% =============================================================
% 第4章 技術成熟度レベルTRLの推移 (課題7) 章本体
% reports.tex から \input される．
% =============================================================

\section{はじめに}
本章には，Arduinoオーケストラ開発における技術成熟度レベル（TRL: Technology Readiness Level）の評価を，
FBS（機能分解構造）の各機能ごとに記録する．
TRLは NASA 由来の評価指標であり一般には9段階で定義されるが，
本プロジェクトは試作開発の初期段階に該当するため，
今回は \textbf{TRL 1〜4 の4段階のみ} で評価する．

\section{TRLレベルの定義（本プロジェクトで用いる範囲）}

\begin{table}[H]
  \centering
  \begin{tabularx}{\linewidth}{cX}
    \toprule
    レベル & 概要 \\
    \midrule
    TRL 1 & 基礎原理が観察・報告された \\
    TRL 2 & 技術コンセプトおよび応用が定式化された \\
    TRL 3 & 解析・実験により概念が実証された \\
    TRL 4 & 構成要素が試験環境で検証された \\
    \bottomrule
  \end{tabularx}
\end{table}

\section{FBS要素ごとのTRL評価}

FBSで定義された各機能要素の現時点におけるTRLレベルと，その評価根拠を表に示す．
FBSの構造はマネジメント計画書（\texttt{management\_plan.tex}）の機能分解構造節を参照．

\begingroup
\small
\begin{longtable}{p{3cm}p{4.5cm}p{1.8cm}p{6cm}}
\toprule
FBSフェーズ & 機能 & TRLレベル & 根拠 \\
\midrule
\endfirsthead

\toprule
FBSフェーズ & 機能 & TRLレベル & 根拠 \\
\midrule
\endhead

\bottomrule
\endfoot

% --- 指揮動作検出 ---
指揮動作検出 & 拍タイミング検出 & TRL2 & 単体テストでの稼働が確認できたため  \\
 & テンポ推定 & TRL3 & テンポ推定単体であれば実測値との誤差が0.025msであり，テンポ変化時も3拍以内の追従を確認できたため \\
 & フェルマータ検出 & TRL1 & 基礎原理（静止検出）は把握済みだが未実装 \\
 & 小節同期見逃し補完 & TRL1 & 基礎原理は把握済みだが未実装 \\
\midrule

% --- 指揮者人形動作 ---
指揮者人形動作 & 腕往復動作 & TRL3 & サーボモータで指揮者棒の往復駆動を実機確認（6月25日）\\
 & 輝度エンコードLED制御 & TRL2 & 腕振りに合わせてLEDを点灯できているため \\
\midrule

% --- 演奏制御 ---
演奏制御 & 輪唱オフセット制御 & TRL3 & 輪唱形式の合奏動作を実機確認（6月25日）\\
 & 楽譜進行制御 & TRL4 & 2分音符・8分音符を含む楽譜制御を試験環境で確認（6月25日）\\
 & テンポ追従演奏 & TRL1 & 基礎原理は把握済みだが未実装 \\
\midrule

% --- 音響合成 ---
音響合成 & 音色合成（倍音・ADSR） & TRL2 & 楽器ごとの音色を出せたため．これ以降は実際の音との波形などと比較する \\
 & エフェクト付与 & TRL1 & 基礎原理は把握済みだが未実装 \\
 & フェルマータ延奏 & TRL1 & 基礎原理は把握済みだが未実装 \\
\midrule

% --- 状態管理 ---
状態管理 & 演奏状態遷移管理 & TRL2 & 予測のみのテストで次の拍予測ができたため \\
 & 異常時フォールバック & TRL1 & 基礎原理は把握済みだが未実装 \\

\end{longtable}
\endgroup
